\section{Rigidity and the coordinate obstruction}\label{sec:rigidity}

Proposition~\ref{prop:lift} converts the hypothetical invariant in $I$ into
a locally nilpotent derivation of $C$ that fixes $v$ and preserves fiber
weights.  We now rule out such a derivation.  An auxiliary grading isolates
the difference of squares $a^2u^2-d^2$ inside $v$.  Its two factors force
three generators to be invariant, leaving a coefficient whose required
fiber weight is impossible.

\begin{proposition}\label{prop:rigidity}
There is no nonzero locally nilpotent derivation of
$C=k[a,d,b,c,u]/(ac-bd-1)$ that preserves the fiber weights
\eqref{bun:fiber-weights} and annihilates
$v=a^3b+a^2u^2-d^2$.
\end{proposition}

\begin{proof}
Suppose that $\widetilde E$ is such a derivation.  Give $C$ the auxiliary
grading in the second row of the following table:
\begin{equation}\label{rig:gradings}
\begin{array}{c|rrrrr}
 &a&d&b&c&u\\ \hline
 \text{fiber weight}&1&1&-1&-1&0\\
 \text{auxiliary degree}&0&1&-1&0&1
\end{array}
\end{equation}
The two gradings are compatible because $ac-bd-1$ is homogeneous of degree
zero for each.  The auxiliary degrees of the two parts of $v$ are
\[
 \deg_{\mathrm{aux}}(a^3b)=-1,\qquad
 \deg_{\mathrm{aux}}(a^2u^2-d^2)=2.
\]

Decompose $\widetilde E$ into its homogeneous components for auxiliary
degree shift, and let $E'$ be the nonzero component with largest shift,
say $m$.  There are only finitely many components: a derivation is
determined by its values on the five homogeneous generators, and each
value is a finite sum of homogeneous terms.  Each component is a
derivation by the product rule.  Moreover, $E'$ is locally nilpotent.
Indeed, for homogeneous $r\in C$ of auxiliary degree $e$, the component
of auxiliary degree $e+nm$ in $\widetilde E^{\,n}(r)$ is $(E')^n(r)$.
For sufficiently large $n$ the former expression vanishes, so the latter
vanishes as well.  Homogeneous elements span $C$, proving local
nilpotence on every element.  Since the gradings in
\eqref{rig:gradings} are compatible, $E'$ still preserves fiber weights.

Taking the component of degree $m+2$ in $\widetilde E(v)=0$ gives
\[
 E'(a^2u^2-d^2)=0.
\]
The factorization
\[
 a^2u^2-d^2=(au-d)(au+d)
\]
and Lemma~\ref{bun:factorial-closure} imply that $E'$ annihilates both
$au-d$ and $au+d$.  Thus $E'(d)=E'(au)=0$, and another application of that
lemma gives
\begin{equation}\label{rig:fixed}
 E'(a)=E'(d)=E'(u)=0.
\end{equation}
All the factors used here are nonzero; for example, after inverting $a$
one has
\begin{equation}\label{rig:a-chart}
 C[a^{-1}]=k[a,a^{-1},d,u][b],\qquad c=\frac{1+bd}{a},
\end{equation}
where their nonvanishing is immediate.

It remains to control the possible values of $E'$ on $b,c$.
Differentiating $ac-bd=1$ and using \eqref{rig:fixed} yields
$aE'(c)=dE'(b)$.  Define
\[
 q=cE'(b)-bE'(c)\in C.
\]
Multiplication by $a$ and $d$, respectively, then gives
\begin{equation}\label{rig:q}
 E'(b)=aq,\qquad E'(c)=dq.
\end{equation}
The element $q$ is nonzero, since otherwise $E'$ would kill every generator.
We claim that $q$ belongs to the polynomial subring $k[a,d,u]$.

The elements $a,d,u$ are algebraically independent by
\eqref{rig:a-chart}.  Put $K=k(a,d,u)$ and localize $C$ at the nonzero
elements of $k[a,d,u]$.  Because these elements lie in $\ker E'$, the
localized derivation remains locally nilpotent.  The resulting algebra is
$K[b]$, with $c=(1+bd)/a$.  A locally nilpotent $K$-derivation of $K[b]$
sends $b$ into $K$: if its value $P(b)$ had degree at least one, iterating
$P(b)\,d/db$ on $b$ would never give zero.  More explicitly, differentiation
of a nonconstant polynomial multiplies its nonzero leading coefficient
by its positive degree, and multiplication by $P$ then leaves the result
nonconstant.  Equation~\eqref{rig:q} therefore gives $q\in K$.

To compute $C\cap K$ inside the fraction field of $C$, use both coordinate
charts:
\[
 C[a^{-1}]=k[a,a^{-1},d,u][b],\qquad
 C[d^{-1}]=k[a,d,d^{-1},u][c].
\]
An element of the first polynomial ring lying in $K$ has no $b$ terms,
and an element of the second lying in $K$ has no $c$ terms.  Hence
\[
 C\cap K\subset
 k[a,a^{-1},d,u]\cap k[a,d,d^{-1},u]=k[a,d,u].
\]
The last equality follows from unique factorization in $k[a,d,u]$:
a reduced denominator that is both a power of $a$ and a power of $d$
is a unit.  The reverse inclusion is immediate, proving the claim.

Finally, preservation of fiber weights makes $E'(b)$ have weight $-1$.
By \eqref{rig:q}, the nonzero polynomial $q$ has weight $-2$.  This is
impossible in $k[a,d,u]$, whose generators have weights $1,1,0$.
The contradiction proves the proposition.
\end{proof}

\begin{proof}[Proof of Theorem~\ref{thm:main}]
By Proposition~\ref{prop:polynomial}, the algebra $A$ is a polynomial ring
in four variables and its element $p$ becomes the polynomial $f$ stated
in the theorem.  Proposition~\ref{prop:stable} makes $p$ a coordinate of
$A[w]$.  If $p$ were a coordinate of $A$,
Proposition~\ref{prop:coordinate-obstruction} would produce a nonzero
locally nilpotent derivation $E$ of $B$ with $I\cap\ker E\ne0$.
Proposition~\ref{prop:lift} would then supply a derivation forbidden by
Proposition~\ref{prop:rigidity}.  Thus $p$ is not a coordinate of $A$.
Transporting these two properties through the polynomial presentation
proves the assertion for $f$.
\end{proof}
